% 自动生成：请勿手动编辑。来源：图片/_gen_exp_figs.py
% 场景01 SL — 子图 b：MIKU
\def\datapath{data/01_crossing_ped/}
\pgfplotsset{
    discard if not/.style 2 args={x filter/.code={
        \edef\tempa{\thisrow{#1}}\edef\tempb{#2}
        \ifx\tempa\tempb\else\def\pgfmathresult{nan}\fi
    }},
    discard if not2/.style n args={4}{x filter/.code={
        \edef\tempa{\thisrow{#1}}\edef\tempb{#2}
        \edef\tempc{\thisrow{#3}}\edef\tempd{#4}
        \ifx\tempa\tempb
            \ifx\tempc\tempd\else\def\pgfmathresult{nan}\fi
        \else\def\pgfmathresult{nan}\fi
    }}
}
\begin{tikzpicture}
\begin{axis}[
    width=0.9\linewidth, height=4.2cm,
    xlabel={$s$ (m)}, ylabel={$l$ (m)},
    xmin=0, xmax=25.0, ymin=-2.2, ymax=2.2,
    grid=both, grid style={gray!25},
    axis line style={thick},
    tick label style={font=\scriptsize},
    label style={font=\footnotesize},
    legend style={at={(0.98,0.02)}, anchor=south east, font=\scriptsize, draw=none, fill=white, fill opacity=0.85},
]
\addplot[fill=bglight, draw=none, forget plot, on layer=axis background] coordinates {(0,-1.875) (25.0,-1.875) (25.0,1.875) (0,1.875)} \closedcycle;

\draw[black!55, semithick] (axis cs:0,-1.875) -- (axis cs:25.0,-1.875);
\draw[black!55, semithick] (axis cs:0,1.875) -- (axis cs:25.0,1.875);
\draw[black!30, dashed, thin] (axis cs:0,0) -- (axis cs:25.0,0);
% 可行带阴影（blocked 时按 blocked_s 截断）
\addplot[name path=lmin_miku, draw=vibrantorange, thin, dashed, forget plot, ]
    table[col sep=comma, x=s, y=l_min, discard if not={mode}{miku}] {\datapath sl.csv};
\addplot[name path=lmax_miku, draw=vibrantorange, thin, dashed, forget plot, ]
    table[col sep=comma, x=s, y=l_max, discard if not={mode}{miku}] {\datapath sl.csv};
\addplot[fill=lightgreen!35, draw=none, forget plot] fill between[of=lmin_miku and lmax_miku];
% 路径（blocked 时按 blocked_s 截断）
\addplot[vibrantgreen, line width=2pt, unbounded coords=discard] table[col sep=comma, x=s, y=l_path, discard if not={mode}{miku}] {\datapath sl.csv};
\addlegendentry{MIKU}


% 障碍物 行人：t=0 起点淡出扫掠 + τ(12) 到达时刻高亮位置
\fill[dangerred, opacity=0.30] (axis cs:12.000,-0.600) circle [radius=4pt];
\fill[dangerred, opacity=0.22] (axis cs:12.000,0.225) circle [radius=4pt];
\fill[dangerred, opacity=0.18] (axis cs:12.000,1.050) circle [radius=4pt];
\draw[->, dangerred, thick, opacity=0.55] (axis cs:12.000,-0.600) -- (axis cs:12.000,1.700);
\node[anchor=west, font=\tiny, color=dangerred, opacity=0.7] at (axis cs:12.15,-0.600) {$t{=}0$};
% τ(12)=1.5s 自车到达时刻行人实际位置 l=+1.275（上侧）
\fill[dangerred, opacity=0.95] (axis cs:12.000,1.275) circle [radius=5pt];
\draw[dangerred, thick] (axis cs:12.000,1.275) circle [radius=6.5pt];
\node[anchor=south west, font=\tiny\bfseries, color=dangerred] at (axis cs:12.15,1.275) {行人@$\tau{=}1.5$s};
% 自车贴下侧窄缝通过
\node[anchor=north, font=\tiny, color=vibrantgreen!60!black] at (axis cs:12.000,-0.72) {自车贴下侧过};
\node[rectangle, fill=deepblue!75, draw=deepblue, thick, minimum width=18pt, minimum height=8pt, inner sep=0pt] (egobox) at (axis cs:0.000,0.000) {};
\fill[white] ([xshift=0pt]egobox.east) -- ([xshift=-3pt,yshift=2.5pt]egobox.east) -- ([xshift=-3pt,yshift=-2.5pt]egobox.east) -- cycle;
\node[anchor=south west, font=\tiny\bfseries, color=deepblue] at (axis cs:0.20,1.10) {EGO};
\end{axis}
\end{tikzpicture}
